%-------------------------
% Resume in LateX
% Author : Sourabh Bajaj
% License : MIT
%------------------------

\documentclass[a4paper,11pt]{article}

\usepackage[T1]{fontenc}
\usepackage{lmodern}
\usepackage[empty]{fullpage}
\usepackage{titlesec}
\usepackage[usenames,dvipsnames]{color}
\usepackage{enumitem}
\usepackage[hidelinks]{hyperref}
\usepackage{fancyhdr}
\usepackage[english]{babel}
\usepackage{tabularx}

\input{glyphtounicode}

\pagestyle{fancy}
\fancyhf{} % Clear all header and footer fields
\fancyfoot{}
\renewcommand{\headrulewidth}{0pt}
\renewcommand{\footrulewidth}{0pt}

% Adjust margins
\addtolength{\oddsidemargin}{-0.5in}
\addtolength{\evensidemargin}{-0.5in}
\addtolength{\textwidth}{1in}
\addtolength{\topmargin}{-.5in}
\addtolength{\textheight}{1.0in}

\urlstyle{same}

\raggedbottom
\raggedright
\setlength{\tabcolsep}{0in}

% Sections formatting
\titleformat{\section}{
  \vspace{-4pt}\scshape\raggedright\large
}{}{0em}{}[\color{black}\titlerule \vspace{-5pt}]

% Ensure that generated PDF is machine readable/ATS parsable
\pdfgentounicode=1

%-------------------------
% Custom commands
\newcommand{\resumeItem}[2]{
  \item\small{
    \textbf{#1}{: #2}
  }
}

\newcommand{\resumeHeading}[4]{
    \begin{tabular*}{0.99\textwidth}[t]{l@{\extracolsep{\fill}}r}
      \textbf{#1} & #2 \\
      \textit{\small #3} & \textit{\small #4} \\
    \end{tabular*}\vspace{-5pt}
}

\newcommand{\resumeSubheading}[4]{
  \vspace{-1pt}\item
    \begin{tabular*}{0.97\textwidth}[t]{l@{\extracolsep{\fill}}r}
      \textbf{#1} & #2 \\
      \textit{\small #3} & \textit{\small #4} \\
    \end{tabular*}\vspace{-5pt}
}

\newcommand{\resumeSubSubheading}[2]{
    \begin{tabular*}{0.97\textwidth}{l@{\extracolsep{\fill}}r}
      \textit{\small #1} & \textit{\small #2} \\
    \end{tabular*}\vspace{-5pt}
}

\newcommand{\resumeItemPlain}[1]{
  \item\small{
    {#1}
  }
}

\newcommand{\resumeSubItem}[2]{\resumeItem{#1}{#2}}

\renewcommand{\labelitemii}{$\circ$}

\newcommand{\resumeSubHeadingListStart}{\begin{itemize}[leftmargin=*]}
\newcommand{\resumeSubHeadingListEnd}{\end{itemize}}
\newcommand{\resumeItemListStart}{\begin{itemize}}
\newcommand{\resumeItemListEnd}{\end{itemize}\vspace{-5pt}}

%-------------------------------------------
%%%%%%  CV STARTS HERE  %%%%%%%%%%%%%%%%%%%%%%%%%%%%


\begin{document}

%----------HEADING-----------------
\begin{tabular*}{\textwidth}{l@{\extracolsep{\fill}}r}
  \textbf{\href{https://github.com/sofia-domracheva}{\Large Sofia Domracheva}} & Email: \href{mailto:firotss@gmail.com}{firotss@gmail.com}\\
  \href{https://github.com/sofia-domracheva}{github.com/sofia-domracheva} & ORCID: \href{https://orcid.org/0009-0002-5343-1667}{0009-0002-5343-1667}\\
\end{tabular*}


%-----------EDUCATION-----------------
\section{Education}
  \resumeSubHeadingListStart
    \resumeSubheading
      {Plovdiv University ``Paisii Hilendarski'' test test}{Plovdiv, Bulgaria}
      {Bachelor's degree in Software Technologies and Design; GPA: 5.60/6.00 (Excellent)}{Sep. 2022 -- July 2026}
  \resumeSubHeadingListEnd


%-----------RESEARCH-----------------
\section{Research Experience}
  \resumeSubHeadingListStart
    \resumeSubheading
      {Plovdiv University ``Paisii Hilendarski''}{Plovdiv, Bulgaria}
      {Undergraduate Researcher, Quantum Machine Learning (B.Sc.\ thesis)}{Feb. 2025 -- Sep. 2026}
      \resumeItemListStart
        \resumeItem{Data Pipeline}
          {Formulated a binary classification task of jet reconstruction quality on simulated ATLAS Open Data (hadronic jets): data access with ROOT/PyROOT, preprocessing in scikit-learn (standardization, PCA, rescaling for angle encoding).}
        \resumeItem{Quantum Kernels}
          {Implemented a hybrid quantum-classical SVM (quantum Gram matrix + classical optimization) with a ZZ feature map in Qiskit and benchmarked it against classical RBF and linear SVMs; evaluated the kernel on a GPU-accelerated simulator (cuQuantum), a noise model from live \texttt{ibm\_marrakesh} calibration, and the real device via Qiskit Runtime (demonstration kernel on \texttt{ibm\_marrakesh} within a 10-minute QPU budget, mean absolute deviation 0.010 vs.\ ideal).}
        \resumeItem{Results}
          {The classical RBF SVM reached ROC-AUC 0.88--0.89 vs.\ at most 0.72 for the QSVM. Measured exponential kernel concentration directly and showed that the gap persists even in noise-free simulation, so better hardware alone would not close it.}
        \resumeItem{Code}
          {\href{https://github.com/sofia-domracheva/quantum-jets}{github.com/sofia-domracheva/quantum-jets}}
      \resumeItemListEnd
  \resumeSubHeadingListEnd


%-----------PUBLICATIONS-----------------
\section{Publications}
  \resumeSubHeadingListStart
    \item\small{\textbf{S.~Domracheva}, A.~Penev, S.~Stavrev, ``Comparative Analysis of Classical and Quantum Support Vector Machines for Large Hadron Collider Data Classification,'' accepted for presentation at the \textit{International Conference Automatics and Informatics (ICAI'26, IEEE)}, Varna, Bulgaria, Oct. 2026.}
  \resumeSubHeadingListEnd


%-----------EXPERIENCE-----------------
\section{Work Experience}
  \resumeSubHeadingListStart

    \resumeSubheading
      {A Data Pro}{Burgas, Bulgaria}
      {Software Engineer}{Sep. 2021 -- Feb. 2026}
      \resumeItemListStart
        \resumeItem{Data Visualization \& Backend}
          {Developed features for a Django/PostgreSQL backend and Python-based visualizations of large datasets.}
        \resumeItem{Infrastructure \& Integrations}
          {Implemented a weekly Celery backup task that archived reports; integrated third-party web services (Google OAuth, reCAPTCHA); worked in a team using internal task-tracking systems.}
      \resumeItemListEnd

  \resumeSubHeadingListEnd


%-----------PROJECTS-----------------
\section{Projects}
  \resumeSubHeadingListStart
    \resumeSubItem{AI-Powered Tabletop RPG Platform (FMI\{Codes\} Hackathon)}
      {Team project: a web-based Dungeons \& Dragons game with an AI Game Master and generated scene images. Responsible for the AI integration: connected the ChatGPT API as the Game Master and self-hosted Stable Diffusion, sending generation requests to its local API.}
    \resumeSubItem{Geospatial Analysis Platform (NOIT Competition)}
      {Engineered the backend architecture for a Django-based platform assisting small businesses in identifying optimal retail locations. Developed data pipelines to process OpenStreetMap data (Overpass API) and GIS shapefiles, calculating regional business density across multiple sectors (pharmacies, restaurants, hotels) using Python and SQLite.}
    \resumeSubItem{Award-Winning News Aggregator Bot (European Researchers' Night 2020)}
      {Developed a Telegram chatbot using Python and the Telebot API to scrape, aggregate, and manage news articles. Implemented web scraping with BeautifulSoup and user data storage using SQLite. Won 2nd place in the ``My Chatbot'' competition in Burgas.}
  \resumeSubHeadingListEnd


%--------SKILLS------------
\section{Programming Skills}
  \resumeSubHeadingListStart
      \resumeSubItem{Languages}{Python, JavaScript, SQL, HTML/CSS}
      \resumeSubItem{Technologies}{Qiskit (Aer, Runtime, Machine Learning), Linux, ROOT/PyROOT, scikit-learn, NumPy, Matplotlib, Django, Celery, PostgreSQL, Git}
  \resumeSubHeadingListEnd

\end{document}
